\documentclass[10pt,a4paper,ragged2e,withhyper]{../core/core}
\input{config.tex}
\input{color.tex}
\input{margin.tex}
\input{commands.tex}

\begin{document}

%% \cvsection{Personal}
\name{Luis Mayta}
\tagline{Arquitecto Cloud | Experto en Kubernetes | AWS, Terraform, Infraestructura como Código}
\photoR{2.5cm}{./assets/images/me}
\personalinfo{%
  % Not all of these are required!
  \email{luismayta4k@gmail.com}
  \phone{(+51) 959-196-850}
  \location{Lima, Peru}
  \homepage{lucho.hadenlabs.com}
  \twitter{slovacus}
  \linkedin{luismayta}
  \github{luismayta}
  \orcid{0000-0003-0759-5687}
}

%% Make the header extend all the way to the right, if you want.
\begin{fullwidth}
\makecvheader
\end{fullwidth}

%-------------------------------------------------------------------------------
%	SECTION TITLE
%-------------------------------------------------------------------------------
\cvsection[page1sidebar]{Experiencia Laboral Reciente}

\cvexperience{Arquitecto Cloud / Líder de Infraestructura Multi-Agente}{Infosis Arquitectura}{Full-time}{Mar 2025 -- Present}{Argentina}

\begin{itemize}
  \item \emphasis{Diseñó y operó infraestructura cloud escalable en AWS soportando orquestación multi-agente con Kubernetes (on-prem).}
  \item \emphasis{Construyó pipelines CI/CD automatizados con GitLab CI/CD usando Python y Go, mejorando velocidad y confiabilidad de despliegues.}
  \item \emphasis{Desarrolló Infraestructura como Código con Terraform para despliegues reproducibles y versionados.}
  \item \emphasis{Containerizó servicios con Docker para garantizar consistencia de entornos en dev, staging y producción.}
  \item \emphasis{Implementó soluciones de observabilidad con Prometheus, Grafana y OpenTelemetry para monitoreo de salud y desempeño del sistema.}
  \item \emphasis{Automatizó tareas operativas y workflows de despliegue con Python y Go, reduciendo intervención manual.}
  \item \emphasis{Aplicó buenas prácticas de seguridad en entornos cloud, incluyendo políticas IAM, comunicación encriptada y control de acceso basado en roles.}
\end{itemize}

\divider

\cvexperience{Lead Software Architect / Líder de Infraestructura Cloud}{Infosis}{Full-time}{Nov 2022 -- Present}{Argentina}

\begin{itemize}
  \item \emphasis{Lideró la transformación de infraestructura cloud, migrando cargas de trabajo a AWS con orquestación Kubernetes on-prem.}
  \item \emphasis{Implementó Infraestructura como Código con Terraform y Terragrunt en múltiples entornos para reproducibilidad y escalabilidad.}
  \item \emphasis{Estableció mejores prácticas cloud-native, patrones de seguridad y estrategias de optimización de costos.}
  \item \emphasis{Diseñó workflows de despliegue automatizados e integró pipelines CI/CD usando Python y Go para acelerar la entrega de software.}
\end{itemize}

\divider

\cvexperience{Senior SRE / Ingeniero DevOps}{Multiple Companies}{Full-time / Consulting}{2017 -- 2022}{LATAM}

\begin{itemize}
  \item \emphasis{Diseñó y operó plataformas cloud-native usando Kubernetes (on-prem), Terraform y AWS a escala.}
  \item \emphasis{Construyó cimientos de observabilidad y pipelines CI/CD usando Python y Go, permitiendo sistemas de producción confiables.}
  \item \emphasis{Automatizó procesos operativos y despliegues con Python y Go, mejorando la fiabilidad del sistema.}
  \item \emphasis{Optimizó costos y rendimiento de infraestructura en múltiples entornos de alto tráfico.}
  \item \emphasis{Aplicó mejores prácticas de seguridad y cumplimiento en toda la infraestructura cloud.}
\end{itemize}

\divider

\clearpage
\end{document}